\documentclass[10pt,a4paper]{article}
\usepackage[utf8]{inputenc}
\usepackage[left=2cm,right=2cm,top=2cm,bottom=2cm]{geometry}
\usepackage{extarrows}
\usepackage{tcolorbox}
\usepackage{siunitx}
\usepackage[thinlines]{easytable}

\title{Formulario Fisica 1}
\date{A.A. 2020/2021}
\author{Enrico Saggiorato}

\newcommand{\deriv}[2][t]{\dfrac{d{#2}}{d{#1}}}
\newcommand{\derivtwo}[2][t]{\dfrac{d^2{#2}}{d{#1}^2}}
\newcommand{\abs}[1]{\lvert {#1} \rvert}
\newcommand{\scinot}[2]{{#1} \times 10^{#2}}
\newcommand{\dint}{\displaystyle\int}

\begin{document}
    \begin{titlepage}
    \maketitle
    \end{titlepage}


      \section{Cinematica}
      \subsection{Moto Rettilineo Uniforme}
      \begin{itemize}
  	\item Velocità media data la velocità istantanea:
  	\[
  	v_m = \dfrac{1}{t-t_0} \, \int_{t_0}^{t} v(t) \,dt
  	\]
  	\item Legge oraria: \( x(t) = x_0 + v\cdot(t-t_0) \xlongrightarrow{ t_0 = 0 \,}  x(t) = x_0 + vt  \)
      \end{itemize}
	
	\vspace{5mm}
	\subsection{Moto Rettilineo Uniformemente Accelerato}
	\begin{itemize}
	\item Accelerazione in basea alla posizione
	\[
	\int_{x_0}^x a \,dx = \dfrac{1}{2}v^2 - \dfrac{1}{2}v_0^2  
	\]
	\item \( v(t) = v_0 + a\cdot(t - t_0) \xlongrightarrow{ t_0 = 0 \,} v(t) = v_0 + at \)
	\item Legge oraria: \( x(t) = x_0 + v_0\cdot(t-t_0) + \dfrac{1}{2}a\cdot (t-t_0)^2 
		\xlongrightarrow{ t_0 = 0 \,} x(t) = x_0 + v_0t+ \dfrac{1}{2}at^2\)
	\item Velocità data la posizione: \( v^2 = v_0^2 +2a\cdot(x-x_0) \)
	\end{itemize}

	\vspace{5mm}
	\subsection{Cinematica 3D}
	\subsubsection{Coordinate Cartesiane}
	\begin{itemize}
	\item Traiettoria: \( \vec{r}(t) = x(t)\vec{i} + y(t)\vec{j} + z(t)\vec{k} \)
	\item Velocità: \( \vec{v}(t) = \deriv{x}\vec{i} + \deriv{y}\vec{j} + \deriv{z}\vec{k} \)
	\item Accelerazione: \( \vec{a}(t) = \derivtwo{x}\vec{i} + \derivtwo{y}\vec{j} + \derivtwo{z}\vec{k} \)
	\end{itemize}
	\subsubsection{Coordinate Cilindriche}
	\begin{itemize}
	\item Traiettoria: \( \vec{r}(t) = r\vec{\lambda} + z(t)\vec{k} \)
	\item Velocità: \( \vec{v}(t) = \dot{r}\vec{\lambda} + \dot{\theta}r\vec{\mu} + \dot{z}\vec{k}\)
	\item Accelerazione: \( \vec{a}(t) = (\ddot{r}-r\dot{\theta}^2)\vec{\lambda} + (2\dot{r}\dot{\theta}+r\ddot{\theta})
		\vec{\mu} + \ddot{z}\vec{k} \)
	\end{itemize}
	\( \vec{v}_R = \dot{r}\vec{\lambda} \rightarrow \text{VELOCITÀ RADIALE} \quad 
		\vec{v}_{\theta} = \dot{\theta}r\vec{\mu}  \rightarrow  \text{VELOCITÀ TRASVERSA} \) \\
	\( \vec{a}_R = (\ddot{r}-r\dot{\theta}^2)\vec{\lambda} \rightarrow \text{ACCELERAZIONE RADIALE} \quad  
		\vec{a}_{\theta} = (2\dot{r}\dot{\theta}+r\ddot{\theta})\vec{\mu} \rightarrow \text{ACCELERAZIONE TRASVERSA}\)
	\vspace{5mm}
	\begin{tcolorbox}
	\large \textbf{Derivate dei versori}
	\vspace{3mm}
	\[ \dot{\vec{\lambda}} = \dot{\theta}(t)\vec{\mu} \]
	\[ \dot{\vec{\mu}} = -\dot{\theta}(t)\vec{\lambda} \]
	\end{tcolorbox}
	\subsubsection{Coordinate Intrinseche}
	\begin{itemize}
	\item Traiettoria: s(t)
	\item Velocità: \( \vec{v} = \dot{s}(t)\vec{\tau}(t) \)
	\item Accelerazione: \( \vec{a} = \ddot{s}(t)\vec{\tau}(t) + \dfrac{\dot{s}^2(t)}{\rho(t)}\vec{n}(t) \)
	\end{itemize}

	% Non metto vspace qui perché è l'inizio della pagina
	\subsection{Moto Armonico}
	\begin{tcolorbox}
	\large \textbf{Definizioni}
	\begin{itemize}
	\item \( A : \text{Ampiezza} \) 
	\item \( \omega : \text{Pulsazione} \) 
	\item \( \phi : \text{Fase iniziale} \) 
	\end{itemize}
	\vspace{3mm}
	\large \textbf{Relazioni utili} (noti \(x_0\) e \(v_0\))
	\begin{itemize}
	\item \( \tan\phi = \dfrac{x_0\omega}{v_0} \)
	\item \( A^2 = x_0^2+\dfrac{v_0^2}{\omega_0} \)
	\end{itemize}
	\end{tcolorbox}
	\begin{itemize}
	\item Posizione: \( x(t) = A\sin(\omega t + \phi) \)
	\item Periodo: T = \( \dfrac{2\pi}{\omega} \) \hspace{5mm} Pulsazione: \( \omega = \dfrac{2\pi}{T} \) 
	\item Frequenza: \( f = \dfrac{1}{T} = \dfrac{\omega}{2\pi} \)
	\item Velocità: \( v(t) = \omega A\cos(\omega t+\phi) \)
	\item Accelerazione: 
		\begin{align*}
		a(t) & = -\omega^2 A\sin(\omega t+\phi) \\
		& = -\omega^2x(t) 
		\end{align*}
	\item Condizione necessaria e sufficiente per il moto armonico: \( \dfrac{d^2x}{dt^2}+\omega^2x = 0 \)
	\end{itemize}

	\vspace{5mm}
	\subsection{Moto Circolare Uniforme}
	\begin{itemize}
	\item Accelerazione centripeta: \( \vec{a}_N = \dfrac{v^2}{R}\vec{u}_N \qquad \vec{a}_N = \omega^2R\vec{u}_N \)  
	\item Velocità angolare: \( \omega = \dfrac{v}{R} \rightarrow v = \omega R \)
	\item Periodo: \( T = \dfrac{2\pi R}{v} \rightarrow v = \omega R \)
	\item Velocità in forma vettoriale: \( \vec{v} = \vec{\omega} \times \vec{r} \)
	\end{itemize}

	\vspace{5mm}
	\subsection{Moto Circolare Uniformemente Accelerato}
	\begin{itemize}
	\item Legge oraria: \( \theta(t) = \theta_0 + \omega_0 t + \dfrac{1}{2}\alpha t^2 \)
	\item Accelerazione in forma vettoriale: \( \vec{a} = \vec{\alpha} \times \vec{r} \, + \,  \vec{\omega} \times \vec{v} \)
	\end{itemize}

	\vspace{5mm}
	\subsection{Moto Parabolico}
	\begin{itemize}
	\item Legge oraria delle componenti:
		\begin{align*}
		x(t) &= x_0 + (v_0\cos\theta)\cdot t \\
		y(t) &= y_0 + (v_0\sin\theta)\cdot t - \dfrac{1}{2}gt^2
		\end{align*}
	\item Traiettoria: \( y(x) = y_0 + \tan\theta\cdot(x-x_0) - \dfrac{g}{2v_0^2\cos^2\theta}(x-x_0)^2 \)
	\item Massima altezza (supponendo \( (x_0=0)\) e \( (y_0=0) \): 
		\[
		t_{max} = \dfrac{v_0\sin\theta}{g} \qquad 
		x_{max} = \dfrac{v_0^2\sin2\theta}{2g} \qquad
		y_{max} = \dfrac{v_0^2\sin^2\theta_0}{2g} 
		\]
	\item Gittata:  \( x_G = \dfrac{v_0^2\sin2\theta}{g} \rightarrow \text{doppio del punto di massima altezza}\)
	\end{itemize}

	
	\vspace{10mm}
	\section{Dinamica del Punto Materiale}
	\subsection{Principi di Newton}
	\begin{enumerate}
	\item Un corpo non soggetto a forze permane nel suo stato di quiete se era in quiete, oppure in stato di moto rettilineo uniforme se era in moto rettilineo uniforme
	\item \( \vec{F}_T = m\vec{a} \)
	\item \( \vec{F}_{12} = -\vec{F}_{21} \)
	\end{enumerate}

	\vspace{5mm}
	\subsection{Principali tipi di forze}
	\begin{itemize}
	\item Forza gravitazionale: \( \abs{\vec{F}_G} = \gamma \dfrac{m_1 m_2}{r^2} \quad \text{con} \, \gamma = \scinot{6.673}{-11} \, \si{\newton\metre\squared\per\kilo\gram\squared} \)
	\item Forza elettrostatica: \( \abs{\vec{F}_E} = \dfrac{1}{4\pi\epsilon_0} \dfrac{q_1 q_2}{r^2} \quad \text{con} \, \dfrac{1}{4\pi\epsilon_0} = \scinot{8.99}{9} \, \si{\newton\metre\squared\per\coulomb\squared}\)
	\item Forza elastica: \( \abs{\vec{F}_E} = K\abs{\Delta x} \) \quad vettorialmente: \( \abs{\vec{F}_E} = -K(l-l_0)\vec{i} \)
	\item Tensione: diretta lungo la retta della corda, verso il centro della corda stessa
	\item Forza normale: perpemdicolare al piano di appoggio
	\item Attrito statico: \( F_{s,max} = \mu_s N \)
	\item Attrito dinamico: \( \vec{F}_d = -\mu_d N \dfrac{\vec{v}}{\abs{\vec{v}}} \)
	\item Forza viscosa: Usata quando un corpo si muove in un fluido
		\[
		\vec{F}_v = -\beta\vec{v} \quad \text{con} \, \beta = \gamma\eta
		\]
	\item Forza centripeta: moltiplicando l'accelerazione centripeta per la massa ottengo \( \vec{F}_N = m\vec{a}_N \)
	\end{itemize}

	\vspace{5mm}
	\subsection{Pendolo Semplice}
	\begin{itemize}
	\item Risultante: \( m\vec{g} + \vec{T}_f = m\vec{a} \) \\
	In componenti: 
		\begin{align*}
		R_T &= -mg\sin\theta = ma_T \\
		R_N &= T_f-mg\cos\theta = ma_N
		\end{align*}
	\item Accelerazione in componenti:
		\begin{align*}
		a_T &= -g\sin\theta = L\derivtwo{\theta} \\
		a_N &= \dfrac{T_f}{m} - g\cos\theta = \dfrac{v^2}{L}
		\end{align*}
	\item Moto armonico del pendolo semplice: \( \derivtwo{\theta} + \dfrac{g}{L}\theta  = 0 \rightarrow \omega^2 = \dfrac{g}{L} \quad \text{e} \quad T = \displaystyle2\pi\sqrt{\dfrac{L}{g}} \)
	\item Posizione sull'arco di circonferenza: \( s(t) = L\theta(t) = L\theta_0\sin(\omega t+\phi) \)
	\item Velocità angolare: \( \deriv{\theta(t)} = \omega\theta_0\cos(\omega t+\phi) \)
	\item Velocità lineare: \( v = L\omega\theta_0\cos(\omega t+\phi) \rightarrow v = \omega s(t) \) 
	\item Tensione del filo: \( T_f = m\left(g\cos\theta(t) + \dfrac{v^2}{L}\right) \)
	\end{itemize}

	\vspace{5mm}
	\subsection{Quantità di Moto e Impulso}
	\begin{itemize}
	\item Quantità di moto: \( \vec{p} = m\vec{v} \)
	\item Generalizzazione del secondo principio di Newton: \( \deriv{\vec{p}} = m\vec{a} \)
	\item Impulso infinitesimo: \( d\vec{J} = \vec{F}(t)dt \)
	\item Impulso totale: 
		\[ \vec{J}(t_0, t) = \int_{t_0}^{t}\vec{F}(t) \, dt \quad  \xlongrightarrow{ F \, \text{costante} \, } \quad \vec{J} = \vec{F}\cdot(t-t_0) \]
	\item Proprietà dell'impulso: 
		\begin{enumerate}
		\item La somma degli impulsi dovuti a più forze è uguale all'impulso della somma delle forze.
		\item \( \vec{J}(t_0, t_1) + \vec{J}(t_1, t_2) =  \vec{J}(t_0, t_2) \) 
		\end{enumerate}
	\item Teorema dell'impulso: \( \vec{J}(t, t_0) = \Delta\vec{p} \) \\
	Da cui deriva :
	\item Principio di conservazione della quantità di moto: \( \vec{F}_tot = 0 \Rightarrow \vec{p} \  \text{si conserva} \)
	\end{itemize}	

	\vspace{5mm}
	\subsection{Energia Cinetica e Lavoro}
	\begin{itemize}
	\item Energia cinetica: \( E_k = \dfrac{1}{2}mv^2 \)
	\item Lavoro infinitesimo: \( dW = \vec{F}(x, y, z)\cdot d\vec{r} \)
	\item Lavoro totale da A a B: \(W_{AB} = \dint_{A, \gamma}^B F_x\,dx+F_y\,dy+F_z\,dz \) \\
		In funzione del tempo:  \(W_{AB} = \dint_{A, \gamma}^B F_x\dot{x}\,dt+F_y\dot{y}\,dt+F_z\dot{z}\,dt \)
	\item Proprietà del lavoro:
		\begin{enumerate}
		\item La somma dei lavori di più forze è uguale al al lavoro della somma delle forze.
		\item \( W_{AB} + W_{BC} = W_{AC} \)
		\end{enumerate}
	\item Lavoro con forza non parallela allo spostamento: \( W = \int_A^B F\cos\theta \, ds \)
		\begin{center}
		\begin{TAB}(r)[5mm]{|c|c|c|}{|c|c|c|}
		\( \theta < \dfrac{\pi}{2} \)  &  \( \cos\theta > 0 \) & Lavoro positivo \\
		\( \theta > \dfrac{\pi}{2} \)  &  \( \cos\theta < 0 \) & Lavoro negativo \\
		\( \theta = \dfrac{\pi}{2} \)  &  \( \cos\theta = 0 \) & Lavoro nullo \\
		\end{TAB}
		\end{center}
	\item Potenza istantanea: \( P = \deriv{L} = F_T\cdot v \), \qquad \(F_T = \text{componente della forza lungo la direzione dello spostamento} \)
	\item Potenza media: \( \dfrac{W}{t} \)
	\item Teorema dell'energia cinetica (o del lavoro): \( W_{AB} = E_k(B) - E_k(A) \)
	\end{itemize}

	

	
	

























\end{document}